\section{The Events of 1858 in Their Nearest Witnesses}

\subsection{The witnesses and their ceilings}

No stenographer stood at Massabielle. What can be known of the
eighteen reported apparitions rests on identifiable strata of
testimony, and this study names the stratum for every claim it
makes. The strata are: (1)~Bernadette Soubirous' own oral accounts
of 1858, given to family, clergy, police, and magistrates, which
survive today chiefly \emph{inside} other documents---interrogation
records, letters, reports---rather than as her own text; (2)~her
later written accounts, composed from 1861 onwards after she had
learned to write; (3)~the contemporaneous administrative and
judicial paper of 1858, whose complicated survival history is
treated at \S2.5; (4)~the 1858--1861 episcopal commission record
summarized in the 1862 mandement (\S4); (5)~early narrative
literature, above all Henri Lasserre's \work{Notre-Dame de Lourdes}
(1869/1870), written with access to living witnesses but as
apologetic narrative, and the much later memoir of Jean-Baptiste
Estrade (1899), an eyewitness of several apparitions writing four
decades after the events; and (6)~the present-day Sanctuary of
Lourdes' institutional chronology, a custodial harmonization of the
earlier strata.\endnote{The stratification is editorial. The
Sanctuary's chronology used throughout this section is ``The
apparitions,''
\url{https://www.lourdes-france.com/en/the-apparitions/} (English)
and \url{https://www.lourdes-france.com/les-apparitions/} (French),
both read 2026-07-25. Lasserre's volume is cited from the 1870
Palm\'e imprint (References). Estrade's \work{Les apparitions de
Lourdes: souvenirs intimes d'un t\'emoin} was not consulted at
first hand (only later translations were located, under copyright);
it is named here as a received witness class, and no claim below
rests on it.} The critical edition of the documentary record is
Ren\'e Laurentin and Bernard Billet, \work{Lourdes: Documents
authentiques} (Paris: Lethielleux, 7 volumes, 1957--1966), which
published the surviving official and private papers of 1858; it was
not accessible to this project, and its existence is recorded here
precisely so that the reader knows where the fuller record lives
(Appendix~A).\endnote{The edition is identified from its standard
bibliographical description; no claim in this study is sourced to
its contents. This is a recorded ceiling, not a verification.}

The Sanctuary's chronology is therefore the working scaffold of
this section---quoted as the custodian's own received account, in
the ``reported'' register, not as a critical edition. Where the
1862 mandement, itself resting on the commission record, states a
fact differently or more cautiously, the mandement is preferred.

\subsection{Lourdes and the Soubirous household}

Bernadette Soubirous was born at the Boly mill in Lourdes on 7
January 1844, the first child of Fran\c{c}ois Soubirous, a miller,
and Louise Cast\'erot, and was baptized two days later. By 1858 the
family had lost the mill and lived in the \term{cachot}, a single
disused room of the former town lock-up. Bernadette, fourteen years
old, small, and asthmatic since a childhood weakened by illness,
could neither read nor write; she had spent the autumn of 1857 at
Bartr\`es minding sheep, returned to Lourdes on 21 January 1858,
and was preparing at last for her First Communion, which she made
on 3 June 1858.\endnote{Sanctuary of Lourdes, ``Bernadette
Soubirous,''
\url{https://www.lourdes-france.com/en/bernadette-soubirous/}
(birth at the Boly mill 7 January 1844; the cachot; her stature,
asthma, illiteracy; Bartr\`es September 1857; return 21 January
1858; First Communion 3 June 1858). The birth and baptism dates and
the parents' names are also given in the canonization
\term{Litterae decretales} of Pius XI, \work{Acta Apostolicae
Sedis} 26 (1934), pp.~73--74, which add that she was entrusted for
some fifteen months to a wet-nurse, Marie Lag\"ues, at Bartr\`es
(\S6).} The poverty was real, but the Sanctuary's own biography
warns against the purely miserabilist portrait: the first ten years
of her childhood were secure and affectionate, and those who
interrogated her in 1858 met, in the words the bishop would adopt,
``clear, precise responses, impressed with a strong
conviction.''\endnote{Sanctuary biography page, quoting the 1862
mandement; the underlying French is at Lasserre 1870, p.~401
(``des r\'eponses nettes, pr\'ecises, pleines d'\`a-propos,
empreintes d'une forte conviction'').}

\subsection{The eighteen reported encounters: a dated inventory}

The table gives the eighteen reported apparitions as the Sanctuary
itself dates and describes them. Three bounds govern the whole
inventory. First, the register is \emph{reported}: the table
transmits the custodial account, not an adjudicated transcript.
Second, the ``fortnight'' asked for on 18 February was not fifteen
consecutive apparitions: the mandement itself records that during
the fortnight, ``when she expected to see every day, she saw
nothing for two days,'' and the Sanctuary's chronology likewise
contains no apparition on 22 or 26 February and reports the morning
of 3 March as empty until a later return.\endnote{Lasserre 1870,
p.~401 (``durant la Quinzaine\ldots\ elle n'a rien vu pendant deux
jours, quoiqu'elle se trouv\^at dans le m\^eme milieu et dans des
circonstances identiques''); Sanctuary chronology for 3 March. The
mandement turns the gaps into an argument against hallucination
(\S4.2).} Third, the crowd figures are the Sanctuary's received
round numbers, useful for scale only.

\begin{apparitionstable}
1 & Thu 11 Feb & Bernadette, her sister, a friend & Gathering wood
by the Gave; ``a noise like a gust of wind''; ``a lady dressed in
white, \ldots a blue belt and a yellow rose on each foot''; sign of
the cross; Rosary; the Lady vanishes.\\
2 & Sun 14 Feb & companions & Bernadette sprinkles holy water at
the vision; ``the lady smiled and bent her head.''\\
3 & Thu 18 Feb & first adults present & First words. Asked to write
her name: ``It is not necessary''; ``I do not promise to make you
happy in this world but in the other''; ``Would you be kind enough
to come here for a fortnight?''\\
4 & Fri 19 Feb & family members & Bernadette brings a lighted
blessed candle---the received origin of the grotto candles.\\
5 & Sat 20 Feb & small crowd & The Lady teaches her a personal
prayer (never published); a great sadness at the end.\\
6 & Sun 21 Feb & about 100 & Early-morning apparition; that day the
police commissioner Jacomet interrogates her; she will call the
vision only \term{Aquer\`o} (``that thing'').\\
7 & Tue 23 Feb & about 150 & A secret ``for her alone.''\\
8 & Wed 24 Feb & crowd & ``Penance! Penance! Penance! Pray to God
for sinners. Kiss the ground as an act of penance for sinners!''\\
9 & Thu 25 Feb & about 300 & ``She told me to go, drink of the
spring\ldots\ I only found a little muddy water. At the fourth
attempt I was able to drink''; the bitter herbs; ``It is for
sinners.''\\
10 & Sat 27 Feb & about 800 & Silent apparition; she drinks at the
spring; penitential gestures.\\
11 & Sun 28 Feb & over 1{,}000 & Ecstasy, prayer, penance; she is
then taken before Judge Ribes, who threatens prison.\\
12 & Mon 1 Mar & over 1{,}500, first priest present & That night
Catherine Latapie plunges her disabled arm into the spring water;
its movement returns (\S5.4).\\
13 & Tue 2 Mar & growing crowd & ``Go and tell the priests that
people are to come here in procession and to build a chapel here'';
Father Peyramale demands the Lady's name and a winter flowering of
the grotto rosebush.\\
14 & Wed 3 Mar & about 3{,}000 (morning) & No vision in the
morning; after school she returns on an inner invitation; asked
again for a name, the vision smiles; Peyramale repeats his demand.\\
15 & Thu 4 Mar & about 8{,}000 & End of the fortnight; the vision
is silent; no name is given. For twenty days Bernadette does not
return, no longer feeling the call.\\
16 & Thu 25 Mar & crowd & The self-identification in Gascon
Occitan, treated fully in \S3; the rosebush does not bloom.\\
17 & Wed 7 Apr & crowd; Dr Douzous observes & The ``miracle of the
candle'': the flame plays over her hand without burning it,
reported at once by the physician Douzous (\S5.5).\\
18 & Fri 16 Jul & from across the Gave & The grotto now fenced off;
from the far bank: ``I felt that I was in front of the Grotto, at
the same distance as before\ldots\ she was more beautiful than
ever!''\\
\end{apparitionstable}

\noindent All quoted phrases are the Sanctuary's English chronology
of the apparitions, read 2026-07-25; the register throughout is the
custodian's received narrative.\endnote{``The apparitions,''
\url{https://www.lourdes-france.com/en/the-apparitions/}. The
French page of the same chronology was read in parallel; its
material differences from the English page concern the 25 March
wording and are treated in \S3.2. The 16 July barrier reflects the
municipal closure of the grotto site in mid-1858; this study did
not verify the administrative acts of closure and reopening at
source and asserts only the barrier as the Sanctuary reports it.}
The eighteen dated encounters of this table are exactly the corpus
the 1862 judgment covers: ``le 11 f\'evrier 1858 et jours suivants,
au nombre de dix-huit fois'' (\S1).

\subsection{The reported words}

The reported message of Lourdes is brief, and its center of gravity
is unmistakable. To Bernadette personally: a promise of happiness
``not in this world but in the other'' (18 February), a personal
prayer and a secret never disclosed (20 and 23 February). To all,
through her: penance, prayer for sinners, and penitential gesture
(24--25 February); procession and a chapel, addressed to ``the
priests'' (2 March); and the name of 25 March (\S3). Nothing in
the received corpus is a doctrine, a prophecy of dated events, a
political instruction, or a new devotion complete with formulas:
the words fold entirely into the Gospel's own call to conversion
and intercession, which is precisely the ground on which later
papal teaching commends Lourdes (\S7).\endnote{The characterization
is editorial synthesis over the Sanctuary chronology quoted above;
Benedict XVI's summary is congruent: the Lady ``instructed her to
deliver certain very simple messages on prayer, penance and
conversion'' (homily at the torchlight procession, Lourdes, 13
September 2008, References).} The private prayer and the secret are
recorded here as received narrative facts about which nothing more
is known; this study reconstructs neither, and no internet
``Lourdes secret'' is admitted into the corpus.

Two received details deserve their exact shape. The Rosary runs
through the whole cycle---Bernadette prays it at the first
apparition and repeatedly thereafter---but the received accounts
describe the Lady joining the doxology rather than reciting the
Aves; Benedict XVI drew from precisely this detail the
``profoundly theocentric character'' of the Rosary.\endnote{Homily
of 13 September 2008: ``it is notable that Bernadette prays the
rosary under the gaze of Mary, who unites herself to her at the
moment of the doxology.''} And the famous fortnight promise---``I
do not promise to make you happy in this world but in the
other''---is a promise to one person, not a general doctrine of
compensation; Pius XII reads it simply as the frame of Bernadette's
subsequent hidden life (\S6).\endnote{Pius XII, \work{Le
p\`elerinage de Lourdes} (2 July 1957), n.~26, English text at
vatican.va (References): Pius XI ``so to say, authenticate[d] on
his part the promise made by the Immaculate to young Bernadette
that she would `be happy not in this world, but in the next.'\,''}

\subsection{The civil authorities: the interrogations at their
ceilings}

The events of 1858 unfolded under the eyes of a skeptical civil
administration, and its officers questioned Bernadette repeatedly:
the police commissioner Jacomet on 21 February, the examining
magistrate (Judge Ribes) on 28 February with a threat of prison,
and the imperial procurator Dutour in the same season; the prefect
Baron Massy pressed the affair from Tarbes.
That the interrogations occurred is multiply witnessed: the
Sanctuary's chronology records Jacomet's interrogation and the
Ribes threat; the Sanctuary's biography of Bernadette preserves the
received sharp exchanges of the Jacomet session (``Aquer\`o---I
didn't say `the Holy Virgin'\ldots\ Monsieur, you've changed it
all''); and Lasserre, writing within a decade, reconstructs the
Jacomet interrogation at length and names Dutour's reports to the
procurator-general.\endnote{Sanctuary chronology (21 and 28
February entries); Sanctuary biography page (the Jacomet
exchanges, quoted as the Sanctuary's received account); Lasserre
1870, pp.~59--76 (the officials and their dispositions; the
interrogation narrative). The officials are identified as Lasserre's narrative
identifies them (``M.~Jacomet''; ``M.~Dutour, Procureur
imp\'erial''; ``le baron Massy''); their forenames circulate in
the later literature and are not asserted here, having not been
re-verified at source.}

The documentary ceiling must be stated with equal care, because it
is itself part of the history. Lasserre records in a footnote of
1869 that he asked the prefecture of the Hautes-Pyr\'en\'ees for
Jacomet's report of the 21 February interrogation and was told the
dossier ``had disappeared,'' whether by disorder, accident, or
interested removal; that the procurator-general's office at Pau
refused on principle to communicate Dutour's reports; and that the
Ministry of Cults refused likewise.\endnote{Lasserre 1870,
footnote printed across pp.~73--74: ``il nous a \'et\'e impossible
d'en avoir communication\ldots\ le dossier relatif \`a cette
affaire avait disparu\ldots'' (prefecture); ``M.~le
Procureur-g\'en\'eral nous a oppos\'e un principe absolu et a
refus\'e de nous communiquer ces pi\`eces''; the Ministry of Cults
``a agi comme le parquet, avec la politesse en moins.''} A century
later the critical edition of Laurentin and Billet assembled and
published the surviving administrative, judicial, and private
documents of 1858; this project could not consult that edition, so
it asserts neither the exact contents of any procès-verbal nor the
completeness of what survives. What the accessible record
establishes is bounded but real: hostile, contemporaneous official
scrutiny existed; it questioned the visionary within days of the
events; and it produced paper whose later fate ran through loss,
refusal, and eventual critical publication. The bishop's
commission worked in the same season with its own interrogations
(\S4), and the mandement's argument leans expressly on how
Bernadette withstood questioning---``never shaken by threats;
to the most generous offers she answered with a noble
disinterestedness.''\endnote{Lasserre 1870, pp.~400--401
(mandement text): ``Soumise \`a de rudes \'epreuves, elle n'a
jamais \'et\'e \'ebranl\'ee par les menaces; aux offres les plus
g\'en\'ereuses, elle a r\'epondu par un noble
d\'esint\'eressement.''}

\subsection{The spring and the first reported cures}

On 25 February, at the ninth apparition, Bernadette scraped at the
muddy ground of the grotto floor and, at the fourth attempt, drank;
a usable spring emerged at the spot and has flowed since. The
received narrative is of uncovering, not of creation from nothing:
neither the Sanctuary nor the 1862 act claims that no subsurface
water existed before that morning, and this study says
``uncovered'' accordingly. What the Sanctuary says about the water
today is deliberately deflationary and is quoted here as the
custodian's own doctrine of the site: the water ``has no
therapeutic properties or specific properties''; ``it is not holy
water''; and it may not be sold. The Sanctuary even preserves
Bernadette's own caution: ``We take the water like medicine\ldots\
We must have faith, we must pray: this water will have no virtue
without faith!''\endnote{Sanctuary of Lourdes, ``The Water of
Lourdes,''
\url{https://www.lourdes-france.com/en/the-water-of-lourdes/},
read 2026-07-25. The chemists' analysis already cited by the 1862
mandement pointed the same way: the cures were worked ``par
l'emploi d'une eau priv\'ee de toute qualit\'e naturelle
curative, au rapport d'habiles chimistes qui en ont fait une
rigoureuse analyse'' (Lasserre 1870, p.~404).}

The first reported cures attached to the spring almost at once. In
the night after the twelfth apparition (1 March 1858), Catherine
Latapie of Loubajac plunged an arm disabled since an accident into
the spring water and recovered its movement; the Sanctuary's own
list of recognized cures opens with her and with six other cures
of 1858, all seven recognized by Bishop Laurence in the mandement
itself on 18 January 1862 (\S4.4, \S5).\endnote{Sanctuary
chronology (1 March); Sanctuary list ``Miraculous healings,''
\url{https://www.lourdes-france.com/en/miraculous-healings/}
(Catherine Latapie through Marie Moreau, each ``Date of
recognition: 18th January 1862'').} From the beginning, then, the
Lourdes pattern was set: the cure is reported at the spring, the
examination is a distinct act, and the recognition belongs to a
bishop. Everything in \S5 is the institutional history of that
distinction.
